\section{Turingmaschinen}

\subsection{Grundidee}
Eine Turing-Maschine (TM) ist ein endlicher Automat, ergänzt um einen Lese-/Schreibkopf und ein beidseitig unendliches Band von Zellen.
Die Eingabe steht zu Beginn auf dem Band; alle übrigen Zellen enthalten das Leerzeichen \(\shortsqcup\).

\subsection{Deterministische Turing-Maschine}
Eine deterministische TM ist ein 7-Tupel
\[
    M = (Q, \Sigma, \Gamma, \delta, q_0, \shortsqcup, F)
\]
mit
\begin{itemize}
    \item endlicher Zustandsmenge \(Q\),
    \item Eingabealphabet \(\Sigma\) mit \(\Sigma \subset \Gamma\),
    \item Bandalphabet \(\Gamma\) (endliche Menge von Symbolen),
    \item Übergangsfunktion \(\delta : Q \times \Gamma \to Q \times \Gamma \times \{L, R\}\),
    \item Startzustand \(q_0 \in Q\),
    \item Leerzeichen \(\shortsqcup \in \Gamma \setminus \Sigma\),
    \item Menge akzeptierender Zustände \(F \subseteq Q\).
\end{itemize}

\subsection{Graphische Darstellung}
Der Übergang \(\delta(q_0, X) = (q_1, Y, D)\) wird dargestellt als\\
\begin{center}
    \begin{tikzpicture}[main/.style = {draw, circle}]
        \node[main] (q) {$q_0$};
        \node[main] (p) [right=1.5cm of q] {$q_1$};
        \draw[->] (q) to node[above] {\(X/Y, D\)} (p);
    \end{tikzpicture}
\end{center}
wobei \(D \in \{L, R\}\) die Bewegungsrichtung des Kopfes angibt.

\subsection{Konfiguration}
Eine Zeichenreihe
\[
    K = X_1 X_2 \ldots X_{i-1}\, q\, X_i X_{i+1} \ldots X_n
\]
heisst Konfiguration von \(M\), wobei \(q \in Q\) der aktuelle Zustand ist und der Kopf über \(X_i\) steht. \\
\textbf{Startkonfiguration:} \(K_0 = q_0 X_1 X_2 \ldots X_n\) (Kopf auf erstem Zeichen).

\subsection{Berechnungsschritte}
Für \(\delta(q, X_i) = (p, Y, R)\) (Bewegung nach rechts):
\[
    \ldots X_{i-1}\, q\, X_i \ldots \vdash \ldots X_{i-1}\, Y\, p\, X_{i+1} \ldots
\]
Für \(\delta(q, X_i) = (p, Y, L)\) (Bewegung nach links):
\[
    \ldots X_{i-1}\, q\, X_i \ldots \vdash \ldots X_{i-2}\, p\, X_{i-1}\, Y\, X_{i+1} \ldots
\]

\subsection{Berechnung}
Eine Berechnung auf Eingabe \(w \in \Sigma^*\) ist eine endliche Folge
\[
    K_0 \vdash K_1 \vdash \cdots \vdash K_n
\]
wobei \(K_0\) die Startkonfiguration ist und aus \(K_n\) keine Bewegung mehr möglich ist. Kurzschreibweise: \(K_0 \vdash^* K_n\).

\subsection{Sprache einer TM}
Die von \(M\) akzeptierte Sprache ist
\[
    L(M) = \{ w \mid q_0 w \vdash^* \alpha p \beta,\; p \in F \text{ und } \alpha, \beta \in \Gamma^* \}.
\]
Eine Sprache, die von einer TM akzeptiert wird, heisst \textbf{rekursiv aufzählbar}. Die Klasse der rekursiv aufzählbaren Sprachen umfasst u.\,a.\ alle kontextfreien Sprachen.

\subsection{Erweiterungen von TM}

\subsubsection{TM mit Speichern}
In der Zustandssteuerung werden zusätzlich endlich viele Datenzustände gespeichert.
Ist nur eine Vereinfachung: \(Q_{\text{neu}} = Q \times \Gamma^k\) für \(k\) Speicher.

\subsubsection{TM mit mehreren Spuren}
Das Band hat \(n\) Spuren; jede Zelle speichert ein \(n\)-Tupel. Simulation durch
\(\Gamma_{\text{neu}} = \Gamma^n\).

\subsubsection{TM mit mehreren Bändern}
Die TM besitzt \(k\) Bänder mit je einem unabhängigen Lese-/Schreibkopf.
Übergangsfunktion:
\[
    \delta : Q \times \Gamma^k \to Q \times \Gamma^k \times \{R, L, S\}^k
\]
\textbf{Satz:} Jede Sprache, die von einer mehrbändigen TM akzeptiert wird, wird auch von einer TM mit nur einem Band akzeptiert (Simulation mit \(2k\) Spuren).

\subsubsection{Nichtdeterministische TM (NTM)}
Eine NTM hat die Übergangsfunktion
\[
    \delta : Q \times \Gamma \to \mathcal{P}(Q \times \Gamma \times D).
\]
\textbf{Satz:} Jede Sprache, die von einer NTM akzeptiert wird, wird auch von einer deterministischen TM akzeptiert (Simulation über Breitensuche aller Konfigurationen).

\subsection{Beschränkungen von TM}

\subsubsection{TM mit semi-unendlichem Band}
Das Band ist nur in eine Richtung unendlich. \\
\textbf{Satz:} Jede TM kann durch eine TM mit semi-unendlichem Band und zwei Spuren simuliert werden.

\subsubsection{Maschine mit mehreren Stacks}
Eine \(k\)-Stack-Maschine ist ein DKA mit \(k\) Stacks.
\[
    \delta : Q \times \Gamma_1 \times \cdots \times \Gamma_k \to Q \times \Gamma_1^* \times \cdots \times \Gamma_k^*
\]
\textbf{Satz:} Jede TM kann durch eine 2-Stack-Maschine simuliert werden (Bewegung durch Verschieben der Inhalte zwischen den beiden Stacks).

\subsubsection{Zählermaschine}
Eine Zählermaschine (Counter Machine) mit \(k\) Zählern entspricht einer \(k\)-Stack-Maschine, bei der Stacks durch Zähler in \(\mathbb{N}\) ersetzt werden. Nur \(= 0\) / \(\neq 0\) unterscheidbar; Inkrement/Dekrement um 1.\\
\textbf{Satz:} Jede TM kann durch eine 2-Zähler-Maschine simuliert werden. (via: TM \(\to\) 2-Stack \(\to\) 3-Zähler \(\to\) 2-Zähler).\\

\textbf{Stack als Zahl (Basis \(k\)):} Der Zustand \(s=X_1X_2\ldots X_n\) eines Stacks mit \(k-1\) Stacksymbolen kann durch einen Zähler \(z\) kodiert werden.
\begin{align*}
    z &= \sum_{i=0}^{n-1} X_{n-i} \cdot k^{n-i-1} \\
    \text{pop}&: z_{\text{new}} = z \div k, \quad X = z \mod k \\
    \text{push}(X_i)&: z_{\text{new}} = z \cdot k + X_i
\end{align*}
Der Dritte Zähler \(k\) dient zur Berechnung.\\

\textbf{3 Zähler \(\to\) 2 Zähler:} Die drei Zähler mit den Werten \(i,j,k\) werden als \(z = 2^i \cdot 3^j \cdot 5^k\) in einem Zähler kodiert. Mit dem zweiten Zähler wird gerechnet.

